\begin{mistake}{2026-07-23}{07}

\begin{problemcontent}
设 \(y_1, y_2\) 是一阶线性非齐次微分方程 \(y' + P(x)y = Q(x)\) 的两个解，若常数 \(\lambda, \mu\)，使得 \(\lambda y_1 + \mu y_2\) 是该方程的解，\(\lambda y_1 - \mu y_2\) 是对应的齐次微分方程的解，则（ ）

A. \(\lambda = -\frac{1}{2}, \mu = -\frac{1}{2}\)

B. \(\lambda = \frac{1}{2}, \mu = \frac{1}{2}\)

C. \(\lambda = \frac{1}{3}, \mu = \frac{2}{3}\)

D. \(\lambda = \frac{2}{3}, \mu = \frac{2}{3}\)
\end{problemcontent}

\begin{solutioncontent}
已知 \(y_1, y_2\) 满足 \(y' + P(x)y = Q(x)\)，将其记为线性算子 \(L[y] = Q(x)\)。

根据题意，\(\lambda y_1 + \mu y_2\) 是非齐次方程的解，代入可得：
\[ L[\lambda y_1 + \mu y_2] = \lambda L[y_1] + \mu L[y_2] = \lambda Q(x) + \mu Q(x) = (\lambda + \mu)Q(x) = Q(x) \]
由此可得：\(\lambda + \mu = 1\)。

又已知 \(\lambda y_1 - \mu y_2\) 是对应齐次方程的解，代入可得：
\[ L[\lambda y_1 - \mu y_2] = \lambda L[y_1] - \mu L[y_2] = \lambda Q(x) - \mu Q(x) = (\lambda - \mu)Q(x) = 0 \]
由此可得：\(\lambda - \mu = 0\)。

联立上述两式，解得：
\[ \lambda = \frac{1}{2}, \mu = \frac{1}{2} \]

故正确选项为 B。
\end{solutioncontent}

\mistakeknowledge{
\begin{itemize}
  \item \textbf{核心公式/定理}：非齐次线性微分方程解的叠加原理。若 \(y_1, y_2\) 是非齐次方程的解，则其线性组合 \(C_1 y_1 + C_2 y_2\) 为非齐次方程解的充要条件是 \(C_1+C_2=1\)；为齐次方程解的充要条件是 \(C_1+C_2=0\)。
  \item \textbf{易错点}：混淆齐次与非齐次方程解的条件，未能直接利用算子代入比较 \(Q(x)\) 的系数来构造方程组。
\end{itemize}}

\end{mistake}
